\section{Two upper-bound mechanisms}\label{sec:proof-overview}

Both proofs use the geometry of the upper spectral edge, but transport
different estimates. The limiting GOE spectral density decays as the
square root of the distance to its upper edge $2$, making $nu^{3/2}$ the
characteristic dimension of an edge window of width $u$. The spectral
estimates below make this count precise at the scales where it is used.
A spherical spin law can be represented as a
Gaussian conditioned on $\|X\|^2=n$, so its resolvent turns this edge
geometry into explicit covariance estimates. The first route observes
all directions at the same precision. The second observes only the soft
edge directions and enlarges that set one spectral cap at a time.

All changes of law below use finite-dimensional likelihood identities
with explicit partition normalizers. The planted law weights the disorder
by its partition function and introduces an auxiliary spin. Estimates
under that law, or under a further tilt, are returned to ordinary disorder
with the stated probability cost. No critical contiguity assertion is
used; the distinction at the spiked-Wigner threshold is discussed in
\cite[Appendix~A]{PerryWeinBandeiraMoitra2018}.

\subsection{Isotropic observations: from covariance to a support profile}

The endpoint of the first route is \Cref{prof:support-profile}: for every
fixed $\omega>0$, on one disorder event of probability tending to one,
\[
 \D_J(f)\ge c_\omega n^{-2/3-\omega}\log(1/p)\,\mu_J(f^2),
 \qquad 0<p:=\mu_J(f\ne0)\le\tfrac12,
\]
simultaneously for every nonzero real $f$. The factor $\log(1/p)$ says
that a function concentrated on a smaller set must pay more heat-bath
energy relative to its squared norm. For an indicator,
$\D_J(\mathbf1_A)=\sum_{x\in A}\mu_J(x)\sum_iQ_i(x,A^c)$ is the
stationary probability current out of $A$ per unit of per-site time.
The profile controls the early evolution of a point-start density;
the spectral gap controls the final decay.

\paragraph{A local covariance comparison.}
At isotropic precision $s$, write $r=\sqrt s$ and $k=nr^3=ns^{3/2}$.
The spectral and spherical estimates in \Cref{prof:sec:spectral} locate
the Gaussian saddle and control its length densities, projections, and
replica overlaps. The unrestricted rotation comparison in
\Cref{prof:sec:unconstrained} then compares the spin and spherical
partition functions and their replica moments. These estimates give the
continuous reference and the probability budgets for the root analysis.

The root analysis studies solutions of
\[
 y+(1-q)m-Jm=h_s,\qquad m=\tanh y,\qquad q=n^{-1}\|m\|^2.
\]
The Kac--Rice formula in \Cref{prof:roots} expresses planted expectations
of sums over regular solutions as an integral against the observation
density and the absolute Jacobian determinant. After gauging the planted
spin to $\mathbf1$, the
conditional matrix calculation in \Cref{prof:matrices} isolates three
directions: $m$; the component of $\mathbf1$ orthogonal to $m$; and the
component of the root field $y$ orthogonal to both. Their span is
$\mathcal P$. The last direction records the remaining cubic part of the
root profile. Conditioning on these directions leaves Haar rotations on
$\mathcal P^\perp$. \Cref{prof:orbits} compares their Haar weights with
the determinant weights in the root count and locates the Gaussian trace.
The exclusions are integrated over the root and conditional matrix data,
with the imposed spectral restriction. The degree-one argument controls
multiplicity when this root count is returned to the planted field law.

Here is the cancellation at the center of the comparison. Put $z=X-m$.
For each replica, fix its projection on $\mathcal P$ and its translated
length; for a pair, also record the overlap of their residual components
in $\mathcal P^\perp$. The two replica types are a product spin prior
with means $m$ and a variable-variance Gaussian prior. Multiplication
by the common quadratic weight
$\exp\{z^{\mathsf T}(J-(1-q)I)z/2\}$ turns the product spin prior
into the posterior $\mu_s$ in the translated coordinates: the root
equation cancels its linear term. The quantities
$\mathcal R_{ij}$ in \eqref{prof:eq16} are their conditional saddle-density
ratios for this residual overlap, relative to a common continuous
reference, for the four ordered type pairs. A projection multiplier first
matches the one-replica saddle densities. The signed combination
$\mathcal R_{pp}-\mathcal R_{pg_0}-\mathcal R_{g_0p}
+\mathcal R_{g_0g_0}$, with the common projection weight in
\eqref{prof:eq16}, cancels common errors and errors depending on only one
replica type. The actual spin statistics retain their lattice measures:
\Cref{prof:common-angular-kernel,prof:mixed-coordinate-transfer} transfer
only the smooth leading terms to common volume, after bounding each
remainder on its own measure.

This signed estimate, including the centered matrix kernel of
\Cref{prof:gaussian-signed-kernel}, yields
\Cref{prof:median-covariance}. For $k\ge\log^{10}n$ and sufficiently small
$r$, most Haar orientations on each retained root orbit satisfy
$\|\Sigma_s\|\le(1+\epsilon)/s$,
together with the weighted covariance and mean bounds in
\eqref{prof:eq13}. Those additional bounds control derivatives along
rotations in the next stage.

\paragraph{A probability bound along the observation path.}
\Cref{prof:prop-amplified-covariance} combines short backward variance
transfers with concentration around the local majority just obtained.
Starting from fixed positive observation times, it bounds, at each
deterministic time with $0<s\le s_0$ and $k\ge K_{\log}\log n$, the
integrated planted mass of
$\mathcal G\cap\{\|\Sigma_s\|>(1+\epsilon)/s\}$ by $e^{-ak}$.
Here $\mathcal G$ is the rotation-invariant good-spectrum event in the
proposition. The positive-time
covariance and terminal gap inputs are proved at the critical parameter
in \Cref{prof:positive-inputs}; the subcritical spectral-gap theorem is
not applied at its endpoint. In the intervening polylogarithmic window,
the induction uses the unrestricted comparison of
\Cref{prof:sec:polylog} in place of root-preserving rotations.

For the bottom range, put $\ell=\log\log(n+1000)$. At each deterministic
time with $k\le K_{\log}\log n$, \Cref{prof:prop-bottom-covariance} bounds
the integrated planted mass of
$\mathcal G\cap\{\|\Sigma_s\|>Cn^{2/3}\ell^C\}$ by $n^{-B}$ for any
fixed $B>0$. Its
finite-replica expansion retains the central prefactors and lattice
covolumes. After conversion to common volume, the alternating sum cancels
every term that omits a replica label; increasing the fixed moment order
then supplies the required probability exponent. These are estimates
along the observation law at deterministic times, with their stated
integrated exceptions.

\paragraph{Entropy and the uniform event.}
\Cref{prof:entropy} first transfers the planted bounds to one typical
ordinary-disorder event, before choosing $f$. After normalizing
$\mu_J(f^2)=1$, the restricted step assumes
$0<p:=\mu_J(f\ne0)\le1/2$ and $\|f\|_\infty^2\le100p^{-2}$.
It treats small, moderate, and order-$n$ values of $\log(1/p)$ separately.
Variance transport uses the bottom-scale estimate in the first range;
replica moments and the sharp covariance bound retain entropy in the
second; square-reweighted mean estimates control the last range. The
product endpoint converts retained entropy to heat-bath energy, giving
\Cref{prof:restricted-profile}. Dyadic layers remove the bounded-test
restriction and prove the support profile stated above.

\subsection{Spectral caps: from a radial deficit to log-Sobolev}

The second route proves \Cref{cap:thm-lsi}: for every fixed $\rho>0$,
with probability tending to one,
\[
 \Ent_{\mu_J}(f^2)\le n^{2/3+\rho}\D_J(f)
 \qquad\text{for every real }f.
\]
This classical logarithmic Sobolev inequality controls the smoothing of
arbitrary densities through hypercontractivity and also implies the
Poincar\'e inequality.

\paragraph{The cap and the radial deficit.}
Let $d=2I-J$. At scale $p$, observe with precision
$B_p=(p^2-d)_+$ and write $J_p=J-B_p$. On the spectral event used below,
the active projection $P_p=\mathbf1_{\{d\le p^2\}}$ has rank of order
$k_p=np^3$; the cap makes $J_p$ scalar on its range. For $h$ in this range, the
spherical comparison uses
$R_p=(\max(d,p^2)+\alpha)^{-1}$, where $\alpha$ is chosen so that
$\operatorname{Tr}R_p+\|R_ph\|^2=n$. \Cref{cap:static-comparison}
controls partition values and projected means uniformly on each prescribed
ball $\|h\|\le R p^2M_p$, where $M_p=\sqrt{np}$ and $R$ is fixed.

Changing the field radially changes the length of the Gaussian mean, so
the sphere constraint also changes $\alpha$. This reduces the radial
response relative to the frozen resolvent $R_p$. The induction records
that reduction explicitly. The true potential is
$\phi_p(h)=\log Z_{J_p}(h)$, whose Hessian is the capped Gibbs covariance.
For this potential,
$K_p=(P_p\nabla^2\phi_p(h)P_p)|_{\operatorname{ran}P_p}$.
During initialization, $K_p$ instead denotes the Hessian of an auxiliary
potential on the same space. Restrict $R_p$ there as well. With
$e=h/\|h\|$, on the prescribed annulus
the matrix $D=p^2(K_p-R_p)$ has a radial entry at most $-0.30$, small
transverse and mixed blocks, and a bounded lower error; see
\eqref{cap:invariant}. This negative radial correction is the
\emph{radial deficit}. It provides the gain needed to propagate the
covariance estimate to a smaller cap.

\paragraph{Transport and initialization.}
The curvature induction proceeds from larger caps to smaller caps.
After the estimates have been proved, the entropy argument uses them
along the forward observation process, in which $p$ and $B_p$ increase.
For a child scale $c=\gamma p$, the exact Gaussian convolution
\eqref{cap:heat-convolution} expresses the child potential through the
parent. The exact increment formula conjugates the covariance of the
parent field under the transition law by the inverse precision increment,
then subtracts that inverse increment, to obtain the child's Hessian.
The Brascamp--Lieb upper bound retains
the radial correction through the inverse in \eqref{cap:bl-recursion};
the variational Fisher lower bound controls the opposite side and the
mixed block. Averaging the rank-one correction regenerates the radial
deficit by \eqref{cap:radial-regeneration}. Meanwhile the still hidden
eigenframe inside the flat parent space is Haar. Its compression makes
the transverse error small enough at the child scale. This conditional
rotation is used under the canonical Gaussian bridge before the change
to the actual tilted transition. \Cref{cap:curvature-induction} is the
resulting one-step statement.

The initial radial deficit comes from the small-field TAP expansion in
\Cref{cap:classical}, specifically \Cref{cap:tap-expansion}.
\Cref{cap:auxiliary-initialization} uses an independent conditional
completion of the first cap to construct an auxiliary potential with the
curvature invariant and a bounded Euclidean gradient error from the true
potential. The companion estimate used for that error is at one fixed
$0<\beta<1$. \Cref{cap:gradient-contraction} contracts the error over a
finite number of cap steps, and \Cref{cap:remove-auxiliary} then transfers
the invariant to the true potential. The induction continues down to
$np^3\asymp\sqrt{\log n}$. For fields $h\in\operatorname{ran}P_p$, its
final output, \Cref{cap:true-curvature}, gives a sharp
$(1+\epsilon)/p^2$ upper bound for the covariance projected onto
$\operatorname{ran}P_p$ on a thin saddle annulus, and a $C/p^2$ bound for
the same projected covariance on any prescribed bounded field ball, at
every scale in the stated interval.

\paragraph{Entropy retention and layers.}
Write $N=\mu_J(f^2)=e^{-L}>0$ and $\nu=f^2\mu_J/N$. The bounded-test
reduction assumes $0\le f\le1$, $\mu_J(f>0)\le1/2$,
$\operatorname{KL}(\nu\Vert\mu_J)\ge L/2$, and $L\le Cn$. The disorder
event is chosen before $f$. For $L\le\lambda n$, where the small fixed
$\lambda>0$ is chosen after the curvature and ordinary-path tolerances,
\Cref{cap:entropy-retention} chooses a first cap with dimension comparable
to $\max\{\sqrt{\log n},c_0L\}$. The curvature bounds control the
information revealed as the cap grows, retaining enough entropy for the
product estimate in \Cref{cap:product-estimate}. For
$\lambda n\le L\le Cn$, \Cref{cap:large-entropy} argues by contradiction
under $\D_J(f)/N\le n^{-2/3}L$. Square-reweighted mean bounds and relative
entropy of stopped observation paths retain order-$n$ entropy until the
product estimate contradicts that assumption. The noise diagnostics and
stability tolerances are chosen before the test function. Finally,
\Cref{cap:layer-lemma} passes from the bounded-test energy estimate to
set capacities, then to arbitrary half-supported functions, and then by
a median decomposition and the entropy shift \eqref{cap:entropy-shift}
to the classical logarithmic Sobolev inequality.

\paragraph{The final dynamical transfer.}
\Cref{sec:dynamics} converts either functional endpoint to worst-start
mixing: the support profile truncates large densities, while logarithmic
Sobolev hypercontractivity supplies the second smoothing argument. The
identity $\mathcal L_J=n(P_J-I)$ and the nonnegative spectrum of $P_J$
keep the factor $n$ between per-site time and attempted updates explicit.
The linear lower tests and \Cref{thm:accepted-critical} supply the
opposite exponents for mixing and the functional constants. The
companion's realized-start and arbitrary-diverging-factor quantifiers
remain those stated in \Cref{thm:accepted-critical}.
